\chapter{The main induction step}
\label{chap:induction}

\section{Restricting to one slice}


\begin{definition}[Appending a line at height $x$]\label{def:append-x}
  For $x \in \F_q$, let $\operatorname{append}_x(\ell)$ be the line in $\F_q^{m+1}$ obtained by appending the coordinate $x$ to every point of a line $\ell \subseteq \F_q^m$. If $f\colon \ell \to \F_q$, let $\operatorname{append}_x(f)$ be the polynomial on $\operatorname{append}_x(\ell)$ defined by $\operatorname{append}_x(f)(u,x)=f(u)$.
\end{definition}

\begin{definition}[$x$-restricted strategy]\label{def:restricted-strategy}
  \uses{def:append-x, def:symmetric-projective-strategy}
  Let $(\psi,A,B,L)$ be a symmetric strategy for the $(m+1,q,d)$ low individual degree test. For each $x \in \F_q$, the $x$-restricted strategy $(\psi,A^x,B^x,L^x)$ for the $(m,q,d)$ test is defined by
  \[
    (A^x)_a^u = A_a^{u,x},
    \qquad
    (B^x)_f^{\ell} = B_{\operatorname{append}_x(f)}^{\operatorname{append}_x(\ell)},
    \qquad
    (L^x)_f^{\ell} = L_{\operatorname{append}_x(f)}^{\operatorname{append}_x(\ell)}.
  \]
\end{definition}

\begin{lemma}[Restricted success probabilities]\label{lem:restricted-probabilities}
  \lean{MIPStarRE.LDT.MainInductionStep.restrictedProbabilities}
  \uses{def:restricted-strategy, def:good-strategy}
  Let $(\psi,A,B,L)$ be an $(\eps,\delta,\gamma)$-good symmetric strategy for the $(m+1,q,d)$ test. For each $x \in \F_q$, let $(\psi,A^x,B^x,L^x)$ be the restricted strategy and let $\eps_x,\delta_x,\gamma_x$ be its three failure probabilities. Then
  \[
    \E_x \eps_x \le \frac{m+1}{m} \eps,
    \qquad
    \E_x \delta_x \le \delta,
    \qquad
    \E_x \gamma_x \le \frac{m+1}{m} \gamma.
  \]
\end{lemma}
\begin{proof}
  The self-consistency test restricts exactly to the slice at height $x$. For the axis-parallel and diagonal-line tests, conditioning on the sampled line not being parallel to the last coordinate identifies the test with the corresponding restricted game.
\end{proof}

\section{The inductive construction}

\begin{theorem}[Self-improvement]\label{thm:self-improvement-in-induction-section}
  \lean{MIPStarRE.LDT.MainInductionStep.selfImprovementInInductionSection}
  \uses{def:polymeasurements, def:simeq, def:strong-self-consistency}
  Let $G \in \polysub{m}{q}{d}$ be a submeasurement such that, on average over $u \sim \F_q^m$,
  \[
    A_a^u \ot I \simeq_\nu I \ot G_{[g(u)=a]}.
  \]
  Set
  \[
    \zeta = 3000m\left(\eps^{1/32}+\delta^{1/32}+(d/q)^{1/32}\right).
  \]
  Then there exists a projective submeasurement $H \in \polysub{m}{q}{d}$ such that:
  \begin{enumerate}
    \item if $H=\sum_h H_h$, then $\bra{\psi} H \ot I \ket{\psi} \ge (1-\nu)-\zeta$;
    \item on average over $u \sim \F_q^m$,
    \[
      A_a^u \ot I \simeq_\zeta I \ot H_{[h(u)=a]};
    \]
    \item $H_h \ot I \approx_\zeta I \ot H_h$;
    \item there exists a positive semidefinite operator $Z$ such that
    \[
      \bra{\psi} Z \ot (I-H) \ket{\psi} \le \zeta,
    \]
    and $Z \ge \E_u A_{h(u)}^u$ for every $h \in \polyfunc{m}{q}{d}$.
  \end{enumerate}
\end{theorem}
\begin{proof}
  \uses{thm:self-improvement, prop:completing-to-measurement, prop:triangle-sub, prop:simeq-data-processing}
  Complete the input submeasurement $G$ by adjoining an extra outcome $\bot$ for its missing mass. Apply Theorem~\ref{thm:self-improvement} to this completed measurement. The output is still indexed by the enlarged answer set, so one then discards the $\bot$ outcome and transfers the estimates back to the original polynomial outcomes. The quantitative work is exactly the completion bookkeeping from Lemma~\ref{prop:completing-to-measurement}: the helper theorem supplies completeness, consistency, strong self-consistency, and boundedness for the completed family, and these pass to the original submeasurement after data processing and one final triangle-inequality step.
\end{proof}

\begin{theorem}[Main induction]\label{thm:main-induction}
  \lean{MIPStarRE.LDT.MainInductionStep.mainInduction}
  \uses{def:good-strategy, def:polymeasurements, def:simeq}
  Let $(\psi,A,B,L)$ be an $(\eps,\delta,\gamma)$-good symmetric strategy for the $(m,q,d)$ low individual degree test. Let $k \ge md$ and set
  \[
    \nu = 1000k^2m^2\left(\eps^{1/1024}+\delta^{1/1024}+\gamma^{1/1024}+(d/q)^{1/1024}\right)
  \]
  and
  \[
    \sigma = m^2\left(\nu + e^{-k/(80000m^2)}\right).
  \]
  Then there exists a measurement $G \in \polymeas{m}{q}{d}$ such that, on average over $u \sim \F_q^m$,
  \[
    A_a^u \ot I \simeq_\sigma I \ot G_{[g(u)=a]}.
  \]
\end{theorem}
\begin{proof}
  \uses{lem:restricted-probabilities, thm:self-improvement-in-induction-section, thm:ld-pasting}
  Proceed by induction on $m$. For $m=1$, the unique axis-parallel line measurement already gives the required global measurement. For the step from $m$ to $m+1$, restrict the strategy to each slice $x$ and apply the induction hypothesis with the same $k$, producing slice measurements $G^x$ with slice-dependent errors $\nu_x$ and $\sigma_x$. Next self-improve each $G^x$ to a projective submeasurement $\widehat G^x$ with a third slice-dependent parameter
  \[
    \zeta_x = 3000m\left(\eps_x^{1/32}+\delta_x^{1/32}+(d/q)^{1/32}\right).
  \]
  Lemma~\ref{lem:restricted-probabilities} and concavity of $\alpha \mapsto \alpha^c$ for $c\le 1$ give averaged bounds $\E_x\nu_x \le \nu$, $\E_x\sigma_x \le \sigma$, and $\E_x\zeta_x \le \zeta$, where here $\zeta$ is the explicit averaged self-improvement error from the source; in particular one shows $\zeta \le \nu$. Theorem~\ref{thm:ld-pasting} then pastes the family $\{\widehat G^x\}$ into a global measurement with error
  \[
    \sigma^* = (\sigma+\zeta)\left(1+\frac{1}{100m}\right) + 2\nu + e^{-k/(80000m^2)}.
  \]
  The final arithmetic check is that $\sigma^* \le (m+1)^2\bigl(\nu+e^{-k/(80000(m+1)^2)}\bigr)$, which is exactly the next-stage value of $\sigma$. This exposes the quantitative cascade hidden behind the one-line induction summary.
\end{proof}

\section{Proof of the main formal theorem}

\begin{proof}
  \proves{thm:main-formal}
  \uses{thm:main-induction, lem:orthonormalization-main-lemma, prop:completing-to-measurement, prop:simeq-triangle-inequality, prop:simeq-data-processing, prop:simeq-to-approx, prop:triangle-sub, lem:schwartz-zippel-individual}
  Start from the original two-prover strategy and symmetrize it by adjoining a classical role register. The symmetrized strategy is still good, now with all three error parameters bounded by $3\eps$, so Theorem~\ref{thm:main-induction} gives a global measurement $G$ on the enlarged space with consistency error $\sigma$. Projecting $G$ to the two role sectors yields POVMs $G^{\mathrm A}$ and $G^{\mathrm B}$ for the original provers. Triangle inequalities together with the original point-measurement consistency show that the post-processed families $G^{\mathrm A}_{[g(u)=a]}$ and $G^{\mathrm B}_{[g(u)=a]}$ are mutually consistent. Schwartz--Zippel then turns this pointwise agreement into global self-consistency, introducing the first intermediate parameter
  \[
    \zeta_1 = 2\sigma + 2\sqrt{3\eps+2\sigma} + \frac{md}{q}.
  \]
  Apply Lemma~\ref{lem:orthonormalization-main-lemma} to obtain projective submeasurements $P^{\mathrm A}$ and $P^{\mathrm B}$, and then complete them to projective measurements $Q^{\mathrm A}$ and $Q^{\mathrm B}$ using Lemma~\ref{prop:completing-to-measurement}. This introduces
  \[
    \zeta_2 = 200\zeta_1^{1/4} + 40\zeta_1^{1/8}.
  \]
  Comparing $Q^{\mathrm A}$ and $Q^{\mathrm B}$ through the original $G^{\mathrm A},G^{\mathrm B}$ gives a second self-consistency parameter
  \[
    \zeta_3 = 6\zeta_1 + 6\zeta_2,
  \]
  and data processing plus one last triangle inequality with the original point measurements gives the final consistency error
  \[
    \zeta_4 = 2\sigma + 2\sqrt{\zeta_1 + \zeta_3/2}.
  \]
  The source then performs the explicit exponent bookkeeping showing that both $\zeta_3/2$ and $\zeta_4$ are bounded by
  \[
    100000k^2m^4\left(\eps^{1/40000} + (d/q)^{1/40000} + e^{-k/(2560000m^2)}\right).
  \]
  Renaming the final projective measurements $Q^{\mathrm A},Q^{\mathrm B}$ as the theorem's $G^{\mathrm A},G^{\mathrm B}$ gives exactly the stated conclusion.
\end{proof}
